\chapter{数据库概念结构设计}

\section{E-R 图符号规范}

本系统采用 Peter Chen E-R 模型作为概念设计工具，图元符号说明如下：

\begin{table}[H]
  \centering
  \caption{E-R 图符号规范}
  \label{tab:er_symbols}
  \begin{tabularx}{\textwidth}{|l|l|X|}
    \hline
    \textbf{图元类型} & \textbf{几何形状} & \textbf{说明} \\
    \hline
    强实体 & 矩形 & 可独立存在、有标识符的事物 \\
    \hline
    弱实体 & 双线矩形 & 依赖其它实体存在的事物 \\
    \hline
    属性 & 椭圆 & 描述实体特征的字段 \\
    \hline
    主键属性 & 下划线椭圆 & 唯一标识实体的属性 \\
    \hline
    联系 & 菱形 & 实体间的关联关系 \\
    \hline
    基数 & 连线标注 & 表示 1:1、1:N、M:N 等关系 \\
    \hline
  \end{tabularx}
\end{table}

\section{局部 E-R 图}

\subsection{局部 E-R 图 1：门店与桌位}

该局部 E-R 图描述门店（Venue）与桌位（Game Table）之间的 1:N 关系。

\begin{figure}[H]
  \centering
  \begin{tikzpicture}[scale=1.1]
    % 实体
    \node[draw, rectangle, thick, fill=blue!20, minimum width=2cm, minimum height=1cm] (venue) at (0,0) {门店};
    \node[draw, rectangle, thick, fill=blue!20, minimum width=2cm, minimum height=1cm] (table) at (6,0) {桌位};

    % 联系
    \node[draw, diamond, thick, fill=orange!20, minimum width=1.5cm, minimum height=1cm] (rel) at (3,0) {包含};

    % 属性
    \node[draw, ellipse, fill=purple!20] (v_id) at (0,2) {id};
    \node[draw, ellipse, fill=purple!20] (v_name) at (-2,1) {name};
    \node[draw, ellipse, fill=purple!20] (v_addr) at (2,1) {address};
    \node[draw, ellipse, fill=purple!20] (t_id) at (6,2) {id};
    \node[draw, ellipse, fill=purple!20] (t_code) at (4,1) {code};
    \node[draw, ellipse, fill=purple!20] (t_pos) at (8,1) {pos\_x, pos\_y};

    % 连线
    \draw (venue) -- node[midway, above] {1} (rel);
    \draw (rel) -- node[midway, above] {N} (table);
    \draw (venue) -- (v_id);
    \draw (venue) -- (v_name);
    \draw (venue) -- (v_addr);
    \draw (table) -- (t_id);
    \draw (table) -- (t_code);
    \draw (table) -- (t_pos);
  \end{tikzpicture}
  \caption{局部 E-R 图 1：门店与桌位}
  \label{fig:er_local1}
\end{figure}

\subsection{局部 E-R 图 2：桌位、预约与会员}

该局部 E-R 图描述桌位（Game Table）、预约（Reservation）和会员（Player）之间的关系。

\begin{figure}[H]
  \centering
  \begin{tikzpicture}[scale=1.0]
    % 实体
    \node[draw, rectangle, thick, fill=blue!20] (table) at (0,0) {桌位};
    \node[draw, rectangle, thick, fill=blue!20] (res) at (4,0) {预约};
    \node[draw, rectangle, thick, fill=blue!20] (player) at (8,0) {会员};

    % 联系
    \node[draw, diamond, thick, fill=orange!20] (rel1) at (2,0) {预订};
    \node[draw, diamond, thick, fill=orange!20] (rel2) at (6,0) {参与};

    % 连线
    \draw (table) -- node[midway, above] {1} (rel1);
    \draw (rel1) -- node[midway, above] {N} (res);
    \draw (res) -- node[midway, above] {0..1} (rel2);
    \draw (rel2) -- node[midway, above] {1} (player);
  \end{tikzpicture}
  \caption{局部 E-R 图 2：桌位、预约与会员}
  \label{fig:er_local2}
\end{figure}

\subsection{局部 E-R 图 3：桌位、对局会话与预约}

\begin{figure}[H]
  \centering
  \begin{tikzpicture}[scale=1.0]
    % 实体
    \node[draw, rectangle, thick, fill=blue!20] (table) at (0,0) {桌位};
    \node[draw, rectangle, thick, fill=blue!20] (session) at (4,0) {对局会话};
    \node[draw, rectangle, thick, fill=blue!20] (res) at (8,0) {预约};

    % 联系
    \node[draw, diamond, thick, fill=orange!20] (rel1) at (2,0) {占用};
    \node[draw, diamond, thick, fill=orange!20] (rel2) at (6,0) {来源};

    % 连线
    \draw (table) -- node[midway, above] {1} (rel1);
    \draw (rel1) -- node[midway, above] {N} (session);
    \draw (session) -- node[midway, above] {0..1} (rel2);
    \draw (rel2) -- node[midway, above] {1} (res);
  \end{tikzpicture}
  \caption{局部 E-R 图 3：桌位、对局会话与预约}
  \label{fig:er_local3}
\end{figure}

\subsection{局部 E-R 图 4：对局、战绩、游戏与会员}

\begin{figure}[H]
  \centering
  \begin{tikzpicture}[scale=0.9]
    % 实体
    \node[draw, rectangle, thick, fill=blue!20] (session) at (0,0) {对局};
    \node[draw, rectangle, thick, fill=blue!20] (record) at (4,0) {战绩};
    \node[draw, rectangle, thick, fill=blue!20] (game) at (8,1) {游戏目录};
    \node[draw, rectangle, thick, fill=blue!20] (player) at (8,-1) {会员};

    % 联系
    \node[draw, diamond, thick, fill=orange!20] (rel1) at (2,0) {产生};
    \node[draw, diamond, thick, fill=orange!20] (rel2) at (6,1) {对应};
    \node[draw, diamond, thick, fill=orange!20] (rel3) at (6,-1) {获胜};

    % 连线
    \draw (session) -- node[midway, above] {1} (rel1);
    \draw (rel1) -- node[midway, above] {N} (record);
    \draw (record) -- node[midway, above] {N} (rel2);
    \draw (rel2) -- node[midway, above] {1} (game);
    \draw (record) -- node[midway, below] {0..1} (rel3);
    \draw (rel3) -- node[midway, below] {1} (player);
  \end{tikzpicture}
  \caption{局部 E-R 图 4：对局、战绩、游戏与会员}
  \label{fig:er_local4}
\end{figure}

\subsection{局部 E-R 图 5：桌位运行态}

桌位运行态是依赖于桌位的弱实体，用于维护当前桌位的实时状态。

\begin{figure}[H]
  \centering
  \begin{tikzpicture}[scale=1.0]
    % 实体
    \node[draw, rectangle, thick, fill=blue!20] (table) at (0,0) {桌位};
    \node[draw, double, rectangle, thick, fill=cyan!20] (state) at (4,0) {桌位状态};

    % 联系
    \node[draw, diamond, thick, fill=orange!20] (rel) at (2,0) {当前状态};

    % 连线
    \draw (table) -- node[midway, above] {1} (rel);
    \draw (rel) -- node[midway, above] {1} (state);
  \end{tikzpicture}
  \caption{局部 E-R 图 5：桌位运行态}
  \label{fig:er_local5}
\end{figure}

\section{全局 E-R 图}

全局 E-R 图整合所有主要实体及其关系，如图 \ref{fig:er_global} 所示。

\begin{figure}[H]
  \centering
  \begin{tikzpicture}[scale=0.8]
    % 实体
    \node[draw, rectangle, thick, fill=blue!20] (venue) at (0,4) {门店};
    \node[draw, rectangle, thick, fill=blue!20] (game) at (6,4) {游戏目录};
    \node[draw, rectangle, thick, fill=blue!20] (table) at (0,2) {桌位};
    \node[draw, double, rectangle, thick, fill=cyan!20] (state) at (0,0) {桌位状态};
    \node[draw, rectangle, thick, fill=blue!20] (player) at (6,2) {会员};
    \node[draw, rectangle, thick, fill=blue!20] (stats) at (8,2) {会员统计};
    \node[draw, rectangle, thick, fill=blue!20] (res) at (2,1) {预约};
    \node[draw, rectangle, thick, fill=blue!20] (session) at (2,-1) {对局会话};
    \node[draw, rectangle, thick, fill=blue!20] (record) at (4,-1) {战绩};

    % 关系连线
    \draw (venue) -- node[midway, left] {1:N} (table);
    \draw (table) -- node[midway, left] {1:1} (state);
    \draw (table) -- node[midway, above] {1:N} (res);
    \draw (player) -- node[midway, right] {1:1} (stats);
    \draw (table) -- node[midway, left] {1:N} (session);
    \draw (res) -- node[midway, left] {0..1:N} (session);
    \draw (session) -- node[midway, above] {1:N} (record);
    \draw (game) -- node[midway, right] {1:N} (record);
    \draw (player) -- node[midway, right] {0..1:N} (record);
  \end{tikzpicture}
  \caption{全局 E-R 图}
  \label{fig:er_global}
\end{figure}

\section{概念设计说明}

从全局 E-R 图可以看出：

\begin{enumerate}
  \item \textbf{门店}与\textbf{桌位}之间是 1:N 关系，一家门店可以有多个桌位
  \item \textbf{桌位状态}是\textbf{桌位}的弱实体，两者是 1:1 关系
  \item \textbf{预约}依赖于\textbf{桌位}，一个桌位可以有多条预约记录
  \item \textbf{对局会话}可以来源于\textbf{预约}，也可以是现场开台（预约为空）
  \item \textbf{战绩}依赖于\textbf{对局会话}和\textbf{游戏目录}，胜者可以是\textbf{会员}或访客
  \item \textbf{会员统计}是\textbf{会员}的一对一扩展表，用于维护聚合数据
\end{enumerate}

该概念模型为后续的逻辑结构设计奠定了坚实基础。
